\chapter{Discussion}\label{chap:discussion}

This chapter interprets the results of Chapter~\ref{chap:results}. It separates what the evidence establishes from what remains uncertain, and it treats the central hypothesis of the project, that an \ac{mpc}-aware choice of learning rule can make secure training cheaper, as still under investigation.

\section{Numerical Agreement between Plain BP and Secure BP}\label{sec:disc-agreement}
For seed 42, the full-size secure run follows the plaintext learning trajectory closely: all five epoch losses and both evaluation losses agree within $4.42\times10^{-4}$, the test counts are identical and the validation counts differ by one example (Table~\ref{tab:res-paired}). This indicates that the batched-matrix program implements the same \ac{bp} computation as the plaintext baseline and that, for this configuration and seed, fixed-point arithmetic with 16 fractional bits did not visibly degrade learning over 200 updates.

Three qualifications apply. First, this is one seed and one run; it does not establish agreement across seeds, and it provides no estimate of run-to-run variation, which probabilistic truncation makes non-zero even for a fixed seed. Second, the secure losses are slightly lower than the plaintext losses in all seven reported values. A systematic sign is not evidence that secure training is better: the differences are about 0.01--0.06\% of the loss values, they arise from different arithmetic (float32 in PyTorch, fixed point with probabilistic truncation and approximate exponentiation, logarithm and division in MP-SPDZ), and a single seed cannot distinguish a systematic bias from chance. Third, the differences are larger than the pre-run fixed-point emulation predicted for the training loss (about $10^{-4}$ at epoch 5 against an observed $2.06\times10^{-4}$), while the observed correct counts (767 and 772) fall within the emulated range (766--767 and 772). The emulation did not model MP-SPDZ's nonlinear approximations, which is a plausible but untested source of the additional deviation. Identical or nearly identical counts do not imply identical predictions or parameters, which the run does not reveal.

\section{Cost of Secure Training}\label{sec:disc-cost}
The full-size secure run took 933.263\,s of MP-SPDZ time and 948.40\,s of external elapsed time, and the three parties together sent 5{,}224.49\,MB. The plaintext training loop for the same configuration took on average 1.467\,s. The two kinds of figure differ by orders of magnitude, but they have different measurement boundaries (the secure figures include preprocessing, the plaintext figure covers only the training loop) and were not shown to come from the same hardware, so no overhead factor is computed. What the figures do establish is that secure \ac{bp} training of this 118{,}282-parameter network on 5{,}000 examples is feasible on one commodity host in about a quarter of an hour under the chosen protocol.

Communication volume is not network latency. All parties ran on one host and communicated over local loopback connections, so the measured time contains no wide-area latency or bandwidth limits. With about $8.9\times10^{6}$ rounds printed for party~0, the run is round-intensive, and in a distributed deployment the number of rounds, not only the volume, would be expected to matter; this is an expectation from the structure of the protocol, not a measurement. The compiler statistics (Table~\ref{tab:impl-scalar-matrix}) show that comparisons for the ReLU masks (about $9.3\times10^{8}$ bit triples), truncations and the per-row softmax and logarithm remain element-wise secure operations after the matrix rewrite. Which of these operations dominates the runtime has not been measured; answering that is the purpose of the profiling step (RQ2) and is required before any \ac{ff} modification can be motivated.

\section{Runtime Measurement Limitations}\label{sec:disc-timing}
The two time measurements answer different questions. The MP-SPDZ time is measured inside party~0 from the start of program execution to its end and, according to the run output, includes the preprocessing generated during the run. The external elapsed time encloses the whole \texttt{compile-run.py} command, which also compiled the program and started and stopped the three party processes. The 15.1\,s difference between the two was not decomposed, and it is not attributed to any particular phase. The timed process tree used 730\,s of user and 1{,}933\,s of system CPU time; the reasons for the high system time were not instrumented. A single run gives no estimate of timing variance, and the host was not dedicated to the experiment under controlled conditions. Future runs should record per-party timings, the host description and, where possible, a separation of online and preprocessing cost.

\section{Matrix versus Scalar Implementation}\label{sec:disc-matrix}
The most consequential implementation decision so far was not a change of learning rule but a change of formulation. The scalar full-size program and the matrix program implement the same \ac{bp} computation, yet their compiled round counts differ by a factor of about 3{,}500 and their truncation counts by a factor of about 72 (Table~\ref{tab:impl-scalar-matrix}); only the matrix program was executable in practice. This has two consequences for the research question. First, any future comparison of secure \ac{bp} and secure \ac{ff} must use the same implementation style, otherwise differences in formulation could outweigh differences in learning rule. Second, part of the potential advantage of \ac{ff} under \ac{mpc}, for example avoiding a backward pass, has to be weighed against an already batched \ac{bp} baseline, not against a naive scalar one.

\section{Fixed-Point Limitations}\label{sec:disc-fixed}
With \texttt{sfix(16,31)}, values are quantised to $2^{-16}\approx1.5\times10^{-5}$ and their magnitude must stay well below $2^{14}$. The single-batch audit of the matrix prototype showed per-element errors of at most $2.2\times10^{-4}$, and very small gradient contributions are rounded to zero or one quantum. Over 200 updates this did not visibly affect the loss trajectory of the seed-42 run, but a smaller learning rate, longer training or larger activations could behave differently. Probabilistic truncation fails with a small probability that grows with the magnitude of the truncated values; a failure would appear as a sudden loss spike, and none is visible in the five epoch losses. This is an observation about one run, not a proof that no truncation error occurred. For \ac{ff}, whose goodness is a sum of squared activities and whose loss and normalisation are nonlinear, range and precision must be analysed again before implementation.

\section{The Truncation-Domain Message}\label{sec:disc-trunc}
The logs of the full-size run, the smoke run, the matrix prototype, A1 and BP-COMP-001 all end with the message that the computation domain is too small for low-round truncation and would need at least 103 bits. The historical A0 log does not contain it; instead, MP-SPDZ suggests enabling \texttt{use\_trunc\_pr} and \texttt{use\_split(3)}, which indicates that A0 was executed without these options (Section~\ref{sec:disc-bounded}). The run completed with exit status 0, and its results agree closely with the plaintext reference, but completion and agreement do not show that the message is without consequence. A reading of the pinned MP-SPDZ source, carried out for this thesis version without executing anything, gives the following picture. The message is printed by the virtual machine's optimisation-suggestion routine after the run. For every probabilistic truncation of a $k$-bit value, the replicated-ring protocol checks whether $k$ plus the statistical security parameter (40 by default) fits into the ring. If it does, a low-round variant is used; if not, a variant with an additional correction step is used, and the largest required size is recorded. The compiled program contains truncations of values with up to 63 bits ($63+40=103$), and with 64-bit shares most truncations in the program cannot use the low-round variant.

On this reading, the message reports that the cheaper truncation variant could not be used, rather than an error. This reading conflicts with an earlier project note, written before the run, which described the message as concerning exact rather than probabilistic truncation; the source reading supersedes that note but has not been reviewed independently. Its consequences are not resolved: the cost of the variant that was executed, its numerical behaviour and its security properties relative to the analysis cited in Section~\ref{sec:bg-fixed} have not been analysed, and whether a larger ring or a different precision would change cost or accuracy has not been tested. The message is therefore recorded as an open technical limitation. It neither invalidates the run nor has been shown to be harmless.

\section{Bounded and Full-Size Secure Runs}\label{sec:disc-bounded}
The bounded runs and the full-size run answer different questions and are kept apart.

\paragraph{Bounded runs.} A0, A1 and BP-COMP-001 show that a small three-party semi-honest training program executes end to end, and for BP-COMP-001 the aggregate outputs agree closely with a matched plaintext FP64 reference. These are feasibility and implementation observations. Their validation and test sets contain only 30 examples each, so one example changes accuracy by $3.33$ percentage points; 13/30 and 14/30 (about 43\% and 47\%) come from a single run and are not evidence of generalisation. The historical runs, at 3/30 (10\%), indicate no useful model quality. BP-COMP-001 differs from A0 and A1 in both example ordering and initialisation, so the improvement cannot be attributed to either change; it is plausible that the class-major schedule and the large-magnitude deterministic initialisation hindered the historical runs, but this has not been tested separately. The bounded cost figures depend on capture coverage and execution conditions and are not compared causally with each other. A re-inspection of the A0 log for this version adds a further difference: the log ends with MP-SPDZ's suggestion to enable probabilistic truncation and the three-way split, and it lacks the truncation-domain message printed by A1 and BP-COMP-001, which indicates that A0 was compiled without the ring-protocol options used by the later runs. This plausibly contributes to its much larger communication (10{,}045.7\,MB against 353.975\,MB), but it has not been verified from the A0 compilation record.

\paragraph{Full-size run.} The limitations of the 30-example sets do not apply to the full-size run, whose 1{,}000-example validation and test sets resolve accuracy in steps of 0.1 percentage points and whose test accuracy of 77.2\% matches the Plain BP run for the same seed. Its cost figures cannot be compared with the bounded runs either: the network, the data size, the program (matrix instead of scalar) and the measurement setup all differ, so no scaling law or speed-up is derived from the two regimes.

\section{Plain BP versus Plain FF}\label{sec:disc-plain}
In the plaintext benchmark, \ac{bp} reached a mean test accuracy of 75.96\% and \ac{ff} 59.76\%, and \ac{bp} trained faster on average (1.467\,s against 3.118\,s for the training loop) while \ac{ff} performed twice as many optimizer updates and four times as many presentations. The comparison is informative about the two implementations in this short-training configuration, but it is not a pure comparison of learning rules: the methods differ in learning rate, objective, update and presentation budget, and inference procedure, and the \ac{ff} configuration was not tuned for this benchmark. Every \ac{ff} seed selected the first epoch of layer~2 as its best validation checkpoint, which suggests that, in this configuration, further layer-2 training did not improve validation accuracy; whether a different threshold, learning rate or negative-sampling scheme changes this was not studied here. The plaintext gap matters for the secure comparison because a cheaper secure \ac{ff} would have to be weighed against its lower utility in this configuration.

\section{Status of the Research Questions and Hypothesis}\label{sec:disc-rq}
RQ1 is only partially addressed: the secure \ac{bp} side has been measured for one seed, but secure \ac{ff} does not exist, so no secure comparison can be made. RQ2 and RQ3 have not been addressed. The central hypothesis, that an \ac{mpc}-aware learning rule can reduce the cost of secure training, is neither supported nor contradicted by the present evidence. Two results bear on how it should be tested: the matrix rewrite shows that implementation style can change secure cost by orders of magnitude for a fixed learning rule, and the plaintext benchmark shows that \ac{ff}, in its current configuration, trains to a markedly lower accuracy than \ac{bp}. Any claim that \ac{ff} reduces \ac{mpc} cost requires a secure \ac{ff} implementation measured under the same configuration.

\section{Limitations and Threats to Validity}\label{sec:disc-limits}
\begin{itemize}
  \item \emph{Statistical coverage.} The full-size secure result is one seed and one run; no confidence statement about secure utility, runtime or communication is possible.
  \item \emph{Baseline provenance.} The seed-42 Plain BP values are repository-recorded evidence from the plaintext benchmark; the plaintext run was not repeated for the paired comparison, and the plaintext hardware is not recorded.
  \item \emph{Evaluation policy.} Plaintext aggregates use validation-selected checkpoints, the secure run uses the final epoch; the two coincide only for seed 42 of Plain BP.
  \item \emph{Measurement scope.} The two time measurements have different boundaries; the external measurement encloses compilation; communication is reported for party~0 and globally, includes preprocessing, and is not latency.
  \item \emph{Environment.} All parties ran on one host; hardware, operating system and network were not recorded in the run artifacts.
  \item \emph{Truncation.} The truncation-domain message is unresolved, the executed truncation variant has not been analysed numerically or for leakage, and no claim of zero leakage is made.
  \item \emph{Threat model.} Results hold only for three parties, a semi-honest adversary and an honest majority under the Rep3 Ring protocol; no malicious security is claimed, and the released losses and counts are not protected by \ac{mpc}.
  \item \emph{Initialisation.} The initial parameters are derived from a public seed; the experiments protect training data, not the initial model.
  \item \emph{Bounded runs.} They use a tiny model and subset and confound ordering with initialisation.
  \item \emph{Scope.} One dataset, small fully connected networks and one protocol; conclusions apply only to this configuration.
\end{itemize}
